\documentclass[11pt]{article}
\usepackage[margin=1in]{geometry}
\usepackage{parskip}
\usepackage[colorlinks=true,allcolors=blue]{hyperref}

\begin{document}

\pagestyle{empty}

\noindent Hikaru Wakaura and Taiki Tanimae\\
QIRI (Quantum Integrated Research Institute Inc.)\\
Tokyo 107-0061, Japan\\
\href{mailto:h.wakaura@qiri.co.jp}{h.wakaura@qiri.co.jp} \quad ORCID: 0000-0001-8381-8323

\vspace{1em}
\noindent July 30, 2026

\vspace{1em}
\noindent To the Editors, \emph{Physical Review B}

\vspace{1em}

\noindent \textbf{Re: Submission of ``Algebraic leak and thermal hybridization,
not athermality, defeat quantum many-body scar codes under generic
measurement''}

\vspace{0.5em}

We submit for your consideration the manuscript above (regular article, two
authors, twelve pages including nine appendices) together with an openly
available simulation package that regenerates every figure with one command.

\vspace{0.5em}

% Paragraph 1: the question and the answer (PRB framing: eigenstate structure)
Quantum many-body scars---non-thermal, weakly entangled, dynamically
coherent eigenstates---are widely regarded as natural candidates for protected
quantum memory. We test this expectation directly in the monitored-dynamics
setting where information protection is sharply defined: encoding a logical
qubit in the canonical PXP scar tower and computing its reference-qubit
coherent information under projective monitoring, benchmarked against an
energy-matched ensemble of thermal encodings. The scar code loses under
every generic measurement we test---local, collective, and coarse-grained
Casimir probes, across the full measurement-rate and block-size grid, with
hierarchical-bootstrap confidence intervals excluding zero at every headline
point---and where information survives, the local deficit \emph{grows} with
system size ($-0.16$ to $-0.38$ for $L=10$ to $18$).

% Paragraph 2: mechanism — the paper's real contribution
The contribution beyond the negative result is mechanistic, and it is
eigenstate-resolved in a way we believe will interest the PRB readership.
Two independent structural liabilities are identified and separated:
\emph{(i)} an algebraic leak---the emergent $\mathrm{su}(2)$ order parameter
that generates the scar tower couples the code states, so any measurement
aligned with it reads the logical qubit out; and \emph{(ii)} thermal
hybridization---local-measurement protection rises monotonically with the
forward-scattering-approximation purity of the rung, crossing over to a
genuine scar \emph{advantage} at the purest rungs. Athermality per se is
neither necessary nor sufficient: what decides the fate of a
weak-ergodicity-breaking code is the interplay of its symmetry algebra with
the measurement basis and its residual overlap with the thermal bulk. An
exactly solvable spin-1 model with zero Casimir variance anchors the clean
limit and correctly predicts the unique fine-tuned probe under which scars
are rescued.

\noindent Relative to prior work: \emph{(i)} the first head-to-head coding
comparison of scar versus thermal subspaces under monitoring at the level of
channel capacity rather than entanglement or fidelity; \emph{(ii)} a
purity-resolved mechanism connecting eigenstate hybridization to operational
code performance; \emph{(iii)} a documented sign-reversing baseline
artifact---single-eigenstate references invert the conclusion; only ensemble
averaging is safe---of direct relevance to anyone benchmarking codes against
individual eigenstates.

\vspace{0.5em}

% Paragraph 3: why PRB
\emph{Physical Review B} is a natural home for this work: quantum many-body
scars, the PXP model, and measurement-induced phase transitions are core PRB
topics, and the paper's central objects---eigenstate hybridization, emergent
symmetry algebras, and their operational consequences---sit squarely in the
condensed-matter tradition of connecting spectral structure to physical
response. The manuscript is fully self-contained (appendices in place of a
supplement), the data are archived with a DOI
(\href{https://doi.org/10.5281/zenodo.21642056}{10.5281/zenodo.21642056}), and
the literature claims have been verified against primary sources.

\vspace{0.5em}

\noindent We suggest the following referees, all leading experts in
measurement-induced dynamics, quantum error correction, or quantum many-body
scars (institutional contact addresses are publicly listed):

\begin{itemize}\itemsep2pt
  \item Prof.\ Matthew P.~A.\ Fisher, University of California, Santa Barbara ---
        measurement-induced transitions
  \item Prof.\ Ehud Altman, University of California, Berkeley ---
        error correction and monitored circuits
  \item Dr.\ Michael J.~Gullans, NIST / University of Maryland ---
        purification transitions
  \item Prof.\ Vedika Khemani, Stanford University --- random and monitored
        circuits
  \item Prof.\ Adam Nahum, \'Ecole Normale Sup\'erieure, Paris --- entanglement
        dynamics
  \item Prof.\ Zlatko Papi\'c, University of Leeds --- quantum many-body scars
  \item Prof.\ Maksym Serbyn, IST Austria --- quantum many-body scars
  \item Dr.\ Sanjay Moudgalya, Technical University of Munich --- exact scars and
        spectrum-generating algebras
\end{itemize}

\noindent We do not request the exclusion of any referee. This manuscript was
previously considered at PRX Quantum, where it was returned at the editorial
stage on selectivity grounds without technical criticism; it has since been
substantially strengthened through three rounds of adversarial internal review
(hierarchical-bootstrap confidence intervals, multi-seed audits, and a
size-dependence analysis added).

\vspace{1em}

\noindent On behalf of both authors,\\
Hikaru Wakaura\\
QIRI (Quantum Integrated Research Institute Inc.), Tokyo 107-0061, Japan\\
\href{mailto:h.wakaura@qiri.co.jp}{h.wakaura@qiri.co.jp}

\end{document}
